% Fichier généré automatiquement par tools/generate_corrections.py.
% Source : DYN/DYN-01/61_Hemostase_02/61_Hemostase_02.tex
% Statut : complet (3/3 question(s) renseignée(s), 0 reprise(s)).
\begin{corrigebox}[Corrigé extrait de la banque]
\CorrigeQuestion{1}
\begin{itemize}
\item On isole la bille. 
\item On réalise le bilan des actions mécaniques :
\begin{itemize}
\item $\torseurstat{T}{\text{rail}}{\text{bille}} = \torseurl{-N_I\vect{z_1}+T_I\vect{x_1}}{\vect{0}}{I}$ $= \torseurl{-N_I\vect{z_1}+T_I\vect{x_1}}{rT_I\vect{y_1}}{G}$.
\item $\torseurstat{T}{\text{bob}}{\text{bille}} = \torseurl{\vect{F}\left(\text{bob} \rightarrow \text{bille} \right)=F(t)\vect{x_0}}{\vect{0}}{G}$.
\item $\torseurstat{T}{\text{fluide}}{\text{bille}} = \torseurl{\vect{F}\left(\text{fluide} \rightarrow \text{bille} \right)=-f_v\vectv{G}{\text{bille}}{0}}{\vect{0}}{G}$.
\item $\torseurstat{T}{g}{\text{bille}} = \torseurl{mg\vect{z_0}}{\vect{0}}{G}$.
\end{itemize}
\item On calcule le torseur dynamique de la bille $\torseurdyn{\text{bille}}{0}=\torseurl{m\vectg{G}{\text{bille}}{0}}{-J\dfrac{R-r}{r}\ddot{\theta}\vect{y_0}}{G}$ avec $\vectg{G}{\text{bille}}{0} = \dfrac{\dd \vectv{G}{\text{bille}}{0}}{\dd t} = (R-r)\ddot{\theta}\vect{x_1} - (R-r)\dot{\theta}^2\vect{z_1}$.
\end{itemize}

En appliquant le TRD à la bille en projection sur $\vect{z_1}$, on a : 
$-N_I+F(t)\vect{x_0}\cdot\vect{z_1} +mg\vect{z_0}\cdot\vect{z_1} = - m(R-r)\dot{\theta}^2$
$\Longleftrightarrow -N_I+F(t)\sin\theta +mg\cos\theta = - m(R-r)\dot{\theta}^2 $.

En appliquant le TMD à la bille, en $G$, en projection sur $\vect{y_0}$, on a : 
$rT_1 = -J\dfrac{R-r}{r}\ddot{\theta}$
$\Longleftrightarrow rT_1 = -\dfrac{2}{5}mr^2 \dfrac{R-r}{r}\ddot{\theta}$
$\Longleftrightarrow T_1 = \dfrac{2}{5}m(r-R)\ddot{\theta}$.
\CorrigeQuestion{2}
En appliquant le TRD à la bille en projection sur $\vect{x_1}$, on a : 
$T_I+F(t)\vect{x_0}\cdot\vect{x_1} +mg\vect{z_0}\cdot\vect{x_1}-f_v \left(R-r\right)\dot{\theta} =  m(R-r)\ddot{\theta}$
$\Longleftrightarrow  T_I+F(t)\cos\theta-mg\sin\theta -f_v \left(R-r\right)\dot{\theta} =  m(R-r)\ddot{\theta}$.

En utilisant la question précédente, on a alors 
$\dfrac{2}{5}m(r-R)\ddot{\theta}+F(t)\cos\theta-mg\sin\theta -f_v \left(R-r\right)\dot{\theta} =  m(R-r)\ddot{\theta}$

$\Longleftrightarrow F(t)\cos\theta = mg\sin\theta  + f_v \left(R-r\right)\dot{\theta} + \dfrac{7}{5}m(R-r)\ddot{\theta}$

(Signe à revoir ?).
\CorrigeQuestion{3}
Si $\theta$ est petit $\dfrac{7}{5}m\left(r-R\right)\ddot{\theta} +f_v\left(r-R\right)\dot{\theta}+mg\theta=F(t)$. 
En passant dans le domaine de Laplace, on a $\dfrac{7}{5}m\left(r-R\right)p^2{\theta(p)} +f_v\left(r-R\right)p {\theta(p)}+mg\theta(p)=F(p)$
soit $H(p)=\dfrac{1}{\dfrac{7}{5}m\left(r-R\right)p^2 +f_v\left(r-R\right)p +mg}$

$ = \dfrac{\dfrac{1}{mg}}{\dfrac{7}{5g}\left(r-R\right)p^2 +\dfrac{f_v}{mg}\left(r-R\right)p +1}$.

On a donc $K_S = \dfrac{1}{mg}$, $\omega_0 = \sqrt{\dfrac{5g}{7\left(r-R\right)}}$,
$\dfrac{2 \xi }{\omega_0} = \dfrac{f_v\left(r-R\right)}{mg}$ soit $\xi =\dfrac{f_v\left(r-R\right)}{2mg}\sqrt{\dfrac{5g}{7\left(r-R\right)}} =\dfrac{f_v}{2mg}\sqrt{\dfrac{5g\left(r-R\right)}{7}} $
\end{corrigebox}
